\setcounter{section}{12}

\inputaufgabe
{}
{

Zeige, dass ein
\definitionsverweis {reelles Vektorbündel}{}{}
über einem Punkt
\zusatzklammer {also einem einpunktigen topologischen Raum} {} {}
das gleiche ist wie ein reeller endlichdimensionaler Vektorraum.

}
{} {}

\inputaufgabe
{}
{

Es sei
\maabb {p} { V } { X
} {}
ein
\definitionsverweis {reelles Vektorbündel}{}{}
über einem
\definitionsverweis {topologischen Raum}{}{}
$X$. Zeige, dass $V$ genau dann ein
\definitionsverweis {Hausdorffraum}{}{}
ist, wenn $X$ ein Hausdorffraum ist.

}
{} {}

\inputaufgabe
{}
{

Es sei $X$ ein
\definitionsverweis {topologischer Raum}{}{.}
Zeige, dass man die
\definitionsverweis {Identität}{}{}
\maabb {\operatorname{Id}_{  X  }} { X } { X
} {}
als ein
\definitionsverweis {reelles Vektorbündel}{}{}
vom
\definitionsverweis {Rang}{}{}
$0$ auffassen kann.

}
{} {}

\inputaufgabe
{}
{

Zeige, dass das Kernbündel aus
[[Reelle Ebene/Parameter/Kernbündel/Beispiel|Beispiel 12.6]]
trivial ist.

}
{} {}

\inputaufgabe
{}
{

Bestimme zum
\definitionsverweis {Vektorbündel}{}{}

\mavergleichskettedisp
{\vergleichskette
{L
}
{ =} { \left\{ (r,s,t,u,v,w)  \mid  ru+sv+tw = 0 , \, (r,s,t) \neq (0,0,0)       \right\}
}
{ \subseteq} { \R^3 \setminus \{0,0,0\} \times \R^3
}
{ \longrightarrow} { \R^3 \setminus \{0,0,0\}
}
{ } {
}
}
{}{}{}
lineare Trivialisierungen oberhalb von
\mathl{D(r)}{,}
\mathl{D(s)}{} und
\mathl{D(t)}{,} also von $r,s,t$ abhängige Basen oberhalb von $D(r)$ usw. Bestimme die Basiswechselabbildungen auf

\mavergleichskette
{\vergleichskette
{D(rs)
}
{ = }{ D(r) \cap D(s)
}
{  }{
}
{    }{
}
{   }{
}
}
{}{}{.}

}
{} {}

\inputaufgabe
{}
{

Bestimme zum
\definitionsverweis {Vektorbündel}{}{}

\mavergleichskettedisp
{\vergleichskette
{L
}
{ =} { \left\{ (r,s,t,u,v,w)  \mid  ru+sv+tw = 0 , \, (r,s,t) \neq (0,0,0)       \right\}
}
{ \subseteq} { \R^3 \setminus \{0,0,0\} \times \R^3
}
{ \longrightarrow} { \R^3 \setminus \{0,0,0\}
}
{ } {
}
}
{}{}{}
diejenigen Parameter
\mathl{(r,s,t)}{,} für die der Vektor $\left(  3   , \,   7  , \,   4                 \right)$ zur Faser
\mathl{L_{(r,s,t)}}{} gehört.

}
{} {}

\inputaufgabe
{}
{

Beschreibe das
\definitionsverweis {Tangentialbündel}{}{}
zur $(n-1 )$-dimensionalen
\definitionsverweis {Sphäre}{}{}

\mavergleichskettedisp
{\vergleichskette
{ S^{n-1}
}
{ =} { \left\{  \left(  x_1   , \,   \ldots  , \,   x_n          \right)   \mid   x_1^2 + \cdots + x_n^2 = 1         \right\}
}
{ \subset} { \R^n
}
{ } {
}
{ } {
}
}
{}{}{}
als ein
\definitionsverweis {Kernbündel}{}{}
im Sinne von
[[Topologischer Raum/Triviale Vektorbündel/Homomorphismus/Kernbündel/Bemerkung|Bemerkung 12.5]],
vergleiche auch
[[Differenzierbare Hyperfläche/Tangentialbündel als Kernbündel/Beispiel|Beispiel 12.7]].

}
{} {}

\inputaufgabe
{}
{

Es sei $X$ ein
\definitionsverweis {topologischer Raum}{}{.}
Zeige, dass ein
\definitionsverweis {Homomorphismus}{}{}
der
\zusatzklammer {trivialen} {} {}
Vektorbündel
\maabbdisp {\varphi} {X \times \R^n } { X \times \R^m
} {}
das gleiche ist wie eine
$m \times n$-\definitionsverweis {Matrix}{}{,}
der Einträge stetige Funktionen von $X$ nach $\R$ sind.

}
{} {}

\inputaufgabe
{}
{

Es sei

\mavergleichskette
{\vergleichskette
{U
}
{ \subseteq }{ \R^n
}
{  }{
}
{    }{
}
{   }{
}
}
{}{}{}
\definitionsverweis {offen}{}{}
und
\maabb {f} { U } { \R^m
} {}
eine
\definitionsverweis {stetig differenzierbare Abbildung}{}{.}
Zeige, dass durch das
\definitionsverweis {totale Differential}{}{}
in der Form
\maabbeledisp {} { U \times \R^n } { U \times \R^m
} { (x,v)} { (x, \left(D f \right)_{ x} (v) )
} {,}
ein Homomorphismus des Vektorbündels
\mathl{U \times \R^n}{} in das Vektorbündel
\mathl{U \times \R^m}{}  gegeben ist.

}
{} {}

\inputaufgabe
{}
{

Betrachte den
\definitionsverweis {topologischen Raum}{}{}

\mavergleichskettedisp
{\vergleichskette
{Y
}
{ \defeq} { \left\{ (s,t,u,v) \in \R^4  \mid  su+tv = 1        \right\}
}
{ } {
}
{ } {
}
{ } {
}
}
{}{}{}
mit der Projektion
\maabbeledisp {p} { Y } { \R^2 \setminus \{ (0,0) \} = X
} {(s,t,u,v)} { (s,t)
} {.}
\aufzaehlungvier{Zeige, dass jede Faser von $p$
\definitionsverweis {homöomorph}{}{}
zu einer reellen Geraden ist.
}{Zeige, dass durch

\mavergleichskettedisp
{\vergleichskette
{ \varphi(s,t)
}
{ =} {(s,t,u(s,t),v(s,t))
}
{ =} { \left( s,t, { \frac{   s  }{ s^2+t^2 } } , { \frac{   t  }{ s^2+t^2 } }  \right)
}
{ } {
}
{ } {
}
}
{}{}{}
eine stetige Abbildung
\maabb {\varphi} { X} { Y
} {}
mit

\mavergleichskettedisp
{\vergleichskette
{ p \circ \varphi
}
{ =} {
\operatorname{Id}_{  X  }
}
{ } {
}
{ } {
}
{ } {
}
}
{}{}{}
gegeben ist.
}{Definiere einen
\definitionsverweis {Homöomorphismus}{}{}
zwischen
\mathkor {} {Y} {und} {X \times \R} {.}
}{Zeige, dass es keine polynomiale Abbildung
\maabb {\psi} { X } { Y
} {}
mit

\mavergleichskettedisp
{\vergleichskette
{ p \circ \psi
}
{ =} {
\operatorname{Id}_{  X  }
}
{ } {
}
{ } {
}
{ } {
}
}
{}{}{}
gibt.
}

}
{} {}

\bild{ \begin{center}
\includegraphics[width=5.5cm]{ \bildeinlesungsvg {Inclusion-exclusion} {svg} }
\end{center}
\bildtext {} }

\bildlizenz { Inclusion-exclusion.svg } {} {Burn~commonswiki} {Commons} {CC-by-sa 3.0} {}

Nicht nur Vektorbündel, sondern allgemeiner topologische Räume kann man durch Verklebungsdaten beschreiben.

Unter einem
\definitionswort {Verklebungsdatum}{}
für
\definitionsverweis {topologische Räume}{}{}
versteht man den folgenden Datensatz.
\aufzaehlungvier{Eine Familie

\mathbed {U_i} {}
 {i \in I} {}
 {} {} {} {,}
von topologischen Räumen.
}{Für jedes Paar
\mathl{(i,j)}{} eine
\definitionsverweis {offene Teilmenge}{}{}

\mavergleichskette
{\vergleichskette
{U_{ij}
}
{ \subseteq }{ U_i
}
{  }{
}
{    }{
}
{   }{
}
}
{}{}{}
\zusatzklammer {mit

\mavergleichskettek
{\vergleichskettek
{U_{ii}
}
{ = }{ U_i
}
{  }{
}
{    }{
}
{   }{
}
}
{}{}{}} {} {.}
}{Für jedes Paar
\mathl{(i,j)}{} einen
\definitionsverweis {Homöomorphismus}{}{}
\maabbdisp {\varphi_{ji}} {U_{ij}} { U_{ji}
} {}
\zusatzklammer {mit

\mavergleichskettek
{\vergleichskettek
{ \varphi_{ii}
}
{ = }{
\operatorname{Id}_{ U_i  }
}
{  }{
}
{    }{
}
{   }{
}
}
{}{}{}} {} {.}
}{Für Indizes

\mavergleichskette
{\vergleichskette
{ i,j,k
}
{ \in }{ I
}
{  }{
}
{    }{
}
{   }{
}
}
{}{}{}
ist die
\definitionswort {Kozykelbedingung}{}

\mavergleichskettedisp
{\vergleichskette
{ \varphi_{kj} \circ \varphi_{ji}
}
{ =} { \varphi_{ki}
}
{ } {
}
{ } {
}
{ } {
}
}
{}{}{}
als Abbildung von
\mathl{U_{ik} \cap U_{ij}}{} nach
\mathl{U_{k}}{} erfüllt.
}

\inputaufgabegibtloesung
{}
{

Es sei ein
\definitionsverweis {Verklebungsdatum}{}{}

\mathbed {U_i} {}
 {i \in I} {}
 {} {} {} {,}
für
\definitionsverweis {topologische Räume}{}{}
gegeben. Zeige, dass es einen eindeutig bestimmten topologischen Raum $X$, eine offene Überdeckung

\mavergleichskette
{\vergleichskette
{ X
}
{ = }{ \bigcup_{i \in I} V_i
}
{  }{
}
{    }{
}
{   }{
}
}
{}{}{}
und
\definitionsverweis {Homöomorphismen}{}{}
\maabb {\psi_i} {U_i } { V_i
} {}
derart gibt, dass

\mavergleichskettedisp
{\vergleichskette
{ \psi_i  \left(  U_{ij}  \right)
}
{ =} { V_i \cap V_j
}
{ } {
}
{ } {
}
{ } {
}
}
{}{}{}
ist und

\mavergleichskettedisp
{\vergleichskette
{ \psi_i \vert_{U_{ij} }
}
{ =} { \psi_j \vert_{U_{ji} }  \circ   \varphi_{ji}
}
{ } {
}
{ } {
}
{ } {
}
}
{}{}{}
gilt.

}
{} {}

\inputaufgabegibtloesung
{}
{

Es sei ein
\definitionsverweis {Verklebungsdatum}{}{}

\mathbed {E_i} {}
 {i \in I} {}
 {} {} {} {,}
über einem
\definitionsverweis {topologischen Raum}{}{}

\mavergleichskettedisp
{\vergleichskette
{ X
}
{ =} { \bigcup_{i \in I} U_i
}
{ } {
}
{ } {
}
{ } {
}
}
{}{}{}
gegeben. Zeige, dass es ein eindeutig bestimmtes reelles Vektorbündel
\maabb {} { E } { X
} {}
und
\definitionsverweis {Isomorphismen}{}{}
\maabb {\psi_i} {E_i } { E \vert_{U_i }
} {}
derart gibt, dass

\mavergleichskettedisp
{\vergleichskette
{ \psi_i \vert_{ E_i {{ |}}_{ U_{i} \cap U_j } }
}
{ =} { \psi_j \vert_{ E_j {{ |}}_{ U_{i} \cap U_j } }  \circ \varphi_{ji}
}
{ } {
}
{ } {
}
{ } {
}
}
{}{}{}
gilt.

}
{} {}

\inputaufgabe
{}
{

Es sei
\maabb {s} { X } { V
} {}
ein
\definitionsverweis {stetiger Schnitt}{}{}
zu einem
\definitionsverweis {reellen Vektorbündel}{}{}
\maabb {p} { V } { X
} {}
über einem
\definitionsverweis {topologischen Raum}{}{}
$X$. Zeige, dass das
\definitionsverweis {Bild}{}{}

\mavergleichskette
{\vergleichskette
{s(X)
}
{ \subseteq }{ V
}
{  }{
}
{    }{
}
{   }{
}
}
{}{}{}
eine
\definitionsverweis {abgeschlossene Teilmenge}{}{}
ist, die
\definitionsverweis {homöomorph}{}{}
zu $X$ ist.

}
{} {}

\inputaufgabe
{}
{

Es sei
\maabb {p} { V } { X
} {}
ein
\definitionsverweis {reelles Vektorbündel}{}{}
vom
\definitionsverweis {Rang}{}{}
$r$ über einem
\definitionsverweis {topologischen Raum}{}{}
$X$ und sei

\mavergleichskette
{\vergleichskette
{ U
}
{ \subseteq }{ X
}
{  }{
}
{    }{
}
{   }{
}
}
{}{}{}
eine
\definitionsverweis {offene Teilmenge}{}{.}
Zeige, dass es auf $U$ eine stetige
\definitionsverweis {Trivialisierung}{}{}
von $V \vert_U$ genau dann gibt, wenn es eine Familie von stetigen
\definitionsverweis {Basisschnitten}{}{}
\maabbdisp {s_1 , \ldots , s_r} { U} {V  \vert_U
} {}
gibt.

}
{} {}

\inputaufgabe
{}
{

Ein
\definitionsverweis {Vektorbündel}{}{}
\maabb {} { V } { X
} {}
vom Rang $m$ über einem
\definitionsverweis {topologischen Raum}{}{}
$X$ sei durch stetige Matrizenabbildungen
\maabbdisp {\varphi_{ij}} { U_i \cap U_j } { \operatorname{GL}_{  m   } \! \left(  \R  \right)
} {}
zu einer offenen Überdeckung

\mavergleichskette
{\vergleichskette
{X
}
{ = }{ \bigcup_{i \in I} U_i
}
{  }{
}
{    }{
}
{   }{
}
}
{}{}{}
gegeben. Zeige, dass ein
\definitionsverweis {stetiger Schnitt}{}{}
\maabbdisp {s} { X } { V
} {}
dasselbe ist wie eine Familie von stetigen Abbildungen
\maabb {t_i} {U_i } { \R^m
} {,}
die die Bedingungen

\mavergleichskettedisp
{\vergleichskette
{t_j
}
{ =} { \varphi_{ij} t_i
}
{ } {
}
{ } {
}
{ } {
}
}
{}{}{}
für alle $ij$ erfüllt.

}
{} {}

\inputaufgabe
{}
{

Es sei

\mavergleichskette
{\vergleichskette
{X
}
{ = }{ \bigcup_{i \in I} U_i
}
{  }{
}
{    }{
}
{   }{
}
}
{}{}{}
eine
\definitionsverweis {offene Überdeckung}{}{}
eines
\definitionsverweis {topologischen Raumes}{}{}
$X$. Es seien
\maabbdisp {\varphi_{ij}} {U_i \cap U_j } { \operatorname{GL}_{  r   } \! \left(  \R  \right)
} {}
und
\maabbdisp {\psi_{ij}} {U_i \cap U_j } { \operatorname{GL}_{  r   } \! \left(  \R  \right)
} {}
\definitionsverweis {Matrixbeschreibungen}{}{,}
die zu den Vektorbündeln
\mathkor {} {E} {bzw.} {F} {}
führen. Zeige, dass diese Bündel genau dann
\definitionsverweis {isomorph}{}{}
sind, wenn es stetige Abbildungen
\maabbdisp {\alpha_i} { U_i } { \operatorname{GL}_{  r   } \! \left(  \R  \right)
} {}
derart gibt, dass
\zusatzklammer {für sinnvolle Einschränkungen} {} {}

\mavergleichskettedisp
{\vergleichskette
{ \alpha_i \circ  \varphi_{ij}  \circ \alpha_j^{-1}
}
{ =} { \psi_{ij}
}
{ } {
}
{ } {
}
{ } {
}
}
{}{}{}
für alle $i,j$ gilt.

}
{} {}

\inputaufgabe
{}
{

Zeige, dass das Komplement des Nullschnittes in einem trivialen
\definitionsverweis {Geradenbündel}{}{}
nicht
\definitionsverweis {zusammenhängend}{}{}
ist.

}
{} {}

\inputaufgabe
{}
{

Man nehme ein schmales rechteckiges Band, verdrehe es mit einer Volldrehung um die längere Achse und verklebe die beiden kürzeren Ränder. Nun schneide man mit einer Schere das Band längs in  der Mitte durch. Ist das entstehende Objekt
\definitionsverweis {zusammenhängend}{}{?}
Wie sieht es aus, wenn man $n$ Halbdrehungen macht?

}
{} {}

Es sei
\maabb {} { V } { X
} {}
ein
\definitionsverweis {reelles Vektorbündel}{}{}
über einem
\definitionsverweis {topologischen Raum}{}{.}
Ein reelles Vektorbündel
\maabb {} { U } { X
} {}
heißt
\definitionswort {Unterbündel}{}
von $V$, wenn ein
\definitionsverweis {injektiver}{}{}
\definitionsverweis {Homomorphismus von Vektorbündeln}{}{}
\maabb {} { U } { V
} {}
gegeben ist.

\inputaufgabe
{}
{

Es sei $X$ ein
\definitionsverweis {topologischer Raum}{}{}
und sei
\maabb {\varphi} { V } { W
} {}
ein
\definitionsverweis {Homomorphismus}{}{}
von
\definitionsverweis {reellen Vektorbündeln}{}{}
über $X$. Es sei

\mavergleichskettedisp
{\vergleichskette
{ Y
}
{ =} { \left\{ P \in X  \mid  \varphi_P \text{ ist surjektiv }         \right\}
}
{ } {
}
{ } {
}
{ } {
}
}
{}{}{}
Zeige, dass die
\zusatzklammer {Familie der} {} {}
\definitionsverweis {Untervektorräume}{}{}

\mavergleichskette
{\vergleichskette
{ \operatorname{kern} \varphi_P
}
{ \subseteq }{ V_P
}
{  }{
}
{    }{
}
{   }{
}
}
{}{}{}
ein
\definitionsverweis {Untervektorbündel}{}{}
von $V\vert_Y$ bilden.

}
{} {}

\inputaufgabe
{}
{

Es sei $X$ ein
\definitionsverweis {topologischer Raum}{}{}
und sei

\mavergleichskette
{\vergleichskette
{ U
}
{ \subseteq }{ V
}
{  }{
}
{    }{
}
{   }{
}
}
{}{}{}
ein
\definitionsverweis {Untervektorbündel}{}{}
des
\definitionsverweis {reellen Vektorbündels}{}{}
\maabb {} { V } { X
} {.}
Zeige, dass die
\zusatzklammer {Familie der} {} {}
\definitionsverweis {Restklassenräume}{}{}

\mathl{V_P/U_P}{} ein Vektorbündel auf $X$ bilden und dass ein
\definitionsverweis {surjektiver}{}{}
\definitionsverweis {Bündelhomomorphismus}{}{}
\maabb {} { V } { V/U
} {}
vorliegt.

}
{} {}

Wir führen zwei weitere Sprechweisen ein.

Es sei $K$ ein
\definitionsverweis {Körper}{}{}
und seien
\mathl{U,V,W}{}
\definitionsverweis {Vektorräume}{}{}
über $K$. Man nennt ein Diagramm der Form

\mathdisp {0 \, \stackrel{  }{\longrightarrow} \, U  \, \stackrel{  }{ \longrightarrow}   \, V  \, \stackrel{  }{\longrightarrow} \, W \,\stackrel{  } {\longrightarrow}  \, 0} {  }

eine \definitionswort {kurze exakte Sequenz}{} von $K$-Vektorräumen, wenn $U$ ein
\definitionsverweis {Untervektorraum}{}{}
von $V$ ist und wenn $W$
\definitionsverweis {isomorph}{}{}
zum
\definitionsverweis {Restklassenraum}{}{}
$V/U$ ist.

Es sei $X$ ein
\definitionsverweis {topologischer Raum}{}{,}
seien
\maabbdisp {} {U,V,W} {X
} {}
\definitionsverweis {Vektorbündel}{}{}
über $X$ und seien
\maabb {\varphi} { U } { V
} {}
und
\maabb {\psi} { V } { W
} {}
\definitionsverweis {Homomorphismen}{}{.}
Man sagt, dass eine
\definitionswort {kurze exakte Sequenz}{}

\mathdisp {0 \, \stackrel{  }{\longrightarrow} \,  U   \, \stackrel{ \varphi }{ \longrightarrow}   \,  V   \, \stackrel{ \psi }{\longrightarrow} \,  W  \,\stackrel{  } {\longrightarrow}  \, 0} {  }

vorliegt, wenn für jeden Punkt

\mavergleichskette
{\vergleichskette
{ P
}
{ \in }{ X
}
{  }{
}
{    }{
}
{   }{
}
}
{}{}{}
für die Fasern eine
\definitionsverweis {kurze exakte Sequenz}{}{}

\mathdisp {0 \, \stackrel{  }{\longrightarrow} \, U_P   \, \stackrel{ \varphi_P }{ \longrightarrow}   \,  V_P   \, \stackrel{ \psi_P }{\longrightarrow} \,  W_P  \,\stackrel{  } {\longrightarrow}  \, 0} {  }

vorliegt.

\inputaufgabe
{}
{

Es sei $X$ ein
\definitionsverweis {topologischer Raum}{}{}
und sei

\mavergleichskette
{\vergleichskette
{ U
}
{ \subseteq }{ V
}
{  }{
}
{    }{
}
{   }{
}
}
{}{}{}
ein
\definitionsverweis {Untervektorbündel}{}{}
des
\definitionsverweis {reellen Vektorbündels}{}{}
\maabb {} { V } { X
} {.}
Zeige, dass eine
\definitionsverweis {kurze exakte Sequenz}{}{}

\mathdisp {0 \, \stackrel{  }{\longrightarrow} \,  U   \, \stackrel{  }{ \longrightarrow}   \,  V   \, \stackrel{  }{\longrightarrow} \,  U/V \,\stackrel{  } {\longrightarrow}  \, 0} {  }

vorliegt.

}
{} {}

\inputaufgabe
{}
{

Es seien
\mathkor {} {X} {und} {Y} {}
\definitionsverweis {topologische Räume}{}{}
und
\maabb {\varphi} { Y } { X
} {}
eine
\definitionsverweis {stetige Abbildung}{}{.}
Es sei
\maabb {p} { V } { X
} {}
ein
\definitionsverweis {Vektorbündel}{}{}
über $X$. Zeige, dass
\mathl{Y \times_X V}{}
\zusatzklammer {siehe
[[Topologische Räume/Relatives Produkt/Aufgabe|Aufgabe 11.16]]} {} {}
ein Vektorbündel über $Y$ ist.

}
{} {}

\inputaufgabe
{}
{

Es seien
\mathkor {} {X} {und} {Y} {}
\definitionsverweis {topologische Räume}{}{}
und
\maabb {\varphi} {Y} {X
} {}
eine
\definitionsverweis {stetige Abbildung}{}{.}
Es sei
\maabb {p} {V} {X
} {}
ein
\definitionsverweis {Vektorbündel}{}{}
über $X$. Beschreibe den
\definitionsverweis {Rückzug}{}{}
eines Vektorbündels mit der Hilfe von
\definitionsverweis {Verklebungsdaten}{}{.}

}
{} {}

\inputaufgabe
{}
{

Es sei $X$ ein
\definitionsverweis {topologischer Raum}{}{}
und seien
\maabb {p} { V } { X
} {}
und
\maabb {q} { W } { X
} {}
\definitionsverweis {Vektorbündel}{}{}
über $X$. Zeige, dass
\mathl{V \times_X W}{}
\zusatzklammer {siehe
[[Topologische Räume/Relatives Produkt/Aufgabe|Aufgabe 11.16]]} {} {}
ein Vektorbündel über $X$ ist, das mit der
\definitionsverweis {direkten Summe von Vektorbündeln}{}{}
über $X$ übereinstimmt.

}
{} {}

\inputaufgabe
{}
{

Es sei
\maabb {\varphi} { X } { Y
} {}
eine
\definitionsverweis {differenzierbare Abbildung}{}{}
zwischen den
\definitionsverweis {differenzierbaren Mannigfaltigkeiten}{}{}
\mathkor {} {X} {und} {Y} {.}
Zeige, dass die zu einem natürlichen
\definitionsverweis {Vektorbündelhomomorphismus}{}{}
\maabbdisp {} {TY} { \varphi^*TX
} {}
von
\definitionsverweis {Vektorbündeln}{}{}
auf $Y$ führt.

}
{} {}

  \zwischenueberschrift{Aufgaben zum Abgeben}

\inputaufgabe
{4}
{

Es sei $X$ ein
\definitionsverweis {topologischer Raum}{}{}
und seien

\mavergleichskette
{\vergleichskette
{ f_1 , \ldots , f_n
}
{ \in }{ C^0(X,\R)
}
{  }{
}
{    }{
}
{   }{
}
}
{}{}{}
reellwertige
\definitionsverweis {stetige Funktionen}{}{}
auf $X$. Es sei

\mavergleichskette
{\vergleichskette
{ U
}
{ = }{ \bigcup_{i \in I} D(f_i)
}
{ \subseteq }{ X
}
{    }{
}
{   }{
}
}
{}{}{}
und

\mavergleichskettedisp
{\vergleichskette
{ S
}
{ =} { \left\{ (P; t_1 , \ldots , f_n)  \mid  P \in U , \,  \sum_{i = 1}^n f_i(P)t_i = 0       \right\}
}
{ \subseteq} { U \times \R^n
}
{ } {
}
{ } {
}
}
{}{}{,}
vergleiche
[[Topologischer Raum/Triviale Vektorbündel/Homomorphismus/Kernbündel/Bemerkung|Bemerkung 12.5]].
Man gebe explizit Trivialisierungen für $S$ über $D(f_i)$ an und zeige, dass $S$ ein Vektorbündel über $U$ vom Rang $n-1$ ist.

}
{} {}

\inputaufgabe
{4}
{

Zeige, dass ein
\definitionsverweis {reelles Geradenbündel}{}{}
\maabb {} { L } { X
} {}
über einem
\definitionsverweis {topologischen Raum}{}{}
$X$ genau dann trivial ist, wenn es einen stetigen nullstellenfreien
\definitionsverweis {Schnitt}{}{}
besitzt.

}
{} {}

\inputaufgabe
{3}
{

Zeige, dass ein Verklebungsdatum für ein Geradenbündel zu einer offenen Überdeckung

\mavergleichskette
{\vergleichskette
{X
}
{ = }{ \bigcup_{i \in I} U_i
}
{  }{
}
{    }{
}
{   }{
}
}
{}{}{}
das gleiche ist wie eine Familie von stetigen nullstellenfreien Funktionen
\maabb {f_{ij}} {U_i \cap U_j } { \R
} {,}
die auf
\mathl{U_i \cap U_j \cap U_k}{} jeweils die Bedingung

\mavergleichskette
{\vergleichskette
{ f_{ki}
}
{ = }{ f_{kj} \cdot f_{ji}
}
{  }{
}
{    }{
}
{   }{
}
}
{}{}{}
erfüllen.

}
{} {}

\inputaufgabe
{4}
{

Zeige, dass das
\definitionsverweis {Möbiusband}{}{}
eine zweidimensionale
\definitionsverweis {Mannigfaltigkeit}{}{}
ist.

}
{} {}

 [[Kategorie:Latexseite]]
